\section{从分销平台到产销平台：Recipe、IP 与信任}

上一章讲推荐与生成，本章继续推演平台形态。王冠认为未来平台会从分销平台走向产销平台：过去生产单点内容，进入公共流量池接受算法匹配；未来生产者可能生产 recipe，即可复用的生产方法，进而形成连续内容和 IP。平台经济的货币也会从注意力转向信任。

\begin{figure}[H]
\centering
\includegraphics[width=0.96\textwidth]{distribution-to-production.png}
\caption{从分销平台到产销平台：单点内容、Recipe、IP 与信任货币。自制概念图，依据 03:20:18--03:42:54 对谈内容整理。}
\end{figure}

\begin{knowledgebox}{读图：稀缺不会消失，只会转移}
分销平台时代，注意力是货币，单点内容进入流量池争抢曝光。产销平台时代，recipe 让内容连续生产，IP 承载稳定风格和质量，信任成为更高价值的货币。
\end{knowledgebox}

\subsection{Recipe：从单点内容到生产方法}

本节把平台形态变化落实到创作者生产的对象上。AI 生成降低单点内容生产成本后，真正有价值的东西会从“这一条内容”转向“持续生产这类内容的方法”。

Recipe 不是某个视频 A 或 B，而是视频 A 和 B 背后的共同生产方法。它可能包含选题、结构、镜头、节奏、风格、素材、目标受众和互动方式。生产 recipe 的人，不只是生产单个内容，而是在生产一套可复用的内容生产能力。

\begin{importantbox}{Recipe 的重要性}
当内容可以低成本生成时，单点内容的稀缺性下降；稀缺性会转移到生产方法、品味、连续性和信任。Recipe 是把这些稀缺性结构化的方式。
\end{importantbox}

\subsection{注意力到信任}

本节讨论平台经济里的“货币”变化。注意力仍然重要，但当内容供给无限扩张时，用户为什么长期跟随某个作者或 IP，会比一次点击更有价值。

传统平台的货币是注意力：播放、停留、点击、互动。王冠认为当内容数量极大增长时，注意力不够分配，价值会转向信任。Substack、Medium、OnlyFans、知识星球、私域社群等都已经在展示这个趋势：用户不是因为平台推荐而消费，而是因为信任作者、频道、IP 或生产方法而付费。

\begin{warningbox}{信任不是注意力的替代，而是更高层货币}
注意力仍然存在，但信任能承载更高客单价、更长期关系和更精准人群。AI 生成让内容供给增加后，可信赖的品味和连续 IP 反而更稀缺。
\end{warningbox}

\subsection{创作者会消失吗}

王冠不认为创作者会消失，而是认为创作者群体会扩大，价值形式会变化。消费本身可能变成创作：用户在使用 GPT-5 或生成系统解决问题时，会产生对他人有价值的内容。创作者不一定亲自剪每一帧、写每一句，而是通过意图、品味、recipe 和信任参与生产。

\subsection{本章小结}

平台从分销走向产销，意味着内容经济的核心资产从单点内容和平台分发，转向 recipe、IP、信任和用户意图。创作者不会简单消失，而会从内容生产者变成生产方法、品味和信任的组织者。

